\documentclass[12pt]{article}
\usepackage[utf8]{inputenc}
\usepackage[T1]{fontenc}
\usepackage{hyperref}
\title{Yoidik Language Resources for K-12 Curriculum Development}
\author{Compiled by Claude Sonnet 5.5 (claude-sonnet-5-5) for Andr\'as Kornai}
\date{}
\begin{document}
\maketitle

\section{Introduction}
Yoidik (ISO 639-3: ydk; Glottocode yoid1240) is a Madang language of the Trans--New Guinea family, spoken in Madang Province, Papua New Guinea. Wikipedia gives about 770 speakers (2003). Joshua Project reports that the Yoidik live along the coast about 50 km north of Madang town and speak Yoidik as their primary language. Matukar people in a neighbouring village use it as a second language. Most Yoidik are Roman Catholic. Search found very little material beyond encyclopedic and mission-listing pages.

\section{Primary Educational Resources}
No primers, textbooks or school materials in Yoidik turned up in search. No evidence of school use was found.

\section{Linguistic and Reference Materials}
\begin{itemize}
\item Wikipedia, Yoidik language (overview, classification, speaker count): \url{https://en.wikipedia.org/wiki/Yoidik_language}
\end{itemize}
No dictionary, grammar or orthography description was found in these results.

\section{Religious and Community Resources}
\begin{itemize}
\item Joshua Project, Yoidik in Papua New Guinea (people-group profile). It lists Scripture status as ``Translation Started'' and says no Bible translation exists in the language yet: \url{https://joshuaproject.net/people_groups/16048}
\item Joshua Project language page for ydk: \url{https://joshuaproject.net/languages/ydk}
\end{itemize}

\section{Digital and Audio Resources}
\begin{itemize}
\item Scripture Earth page for ydk. The search summary mentions a Yoidik page with Global Recordings Network audio, but this was not verified: \url{https://www.scriptureearth.org/00eng.php?iso=ydk}
\end{itemize}

\section{Curriculum Development Context}
Yoidik is a small, apparently stable community language. No published teaching materials, orthography or dictionary were found. Any K-12 curriculum would need first-hand documentation work with the community, and could draw on the language's presumed use as a first language by the whole community. Papua New Guinea's vernacular-preschool and elementary programmes may be relevant, but nothing specific to Yoidik was found.
\end{document}